% =====================================================================
\lezione{12/10/2026}{Recurrent neural network e LSTM}
  {Prof.~Danilo Giordano}
  {perché servono le RNN; hidden state; varianti input/output; RNN mono e bidirezionale; cella RNN e
   classification head; relazione di ricorrenza; BPTT; vanishing ed exploding gradient; dipendenze a
   lungo termine; LSTM; dati in forma di tensore 3D; progettazione di una RNN; da tabella a sequenze}
\label{lez:2026-10-12}
% =====================================================================

\begin{nota}[Fonte]
Trascrizione automatica della registrazione, revisionata. Formule, tabelle e best practice sono
riprese dalle slide \emph{3-deep-learning-methodologies} (pp.~34--71) e dal notebook
\emph{RNN-Pollution-Multivariate}.
\end{nota}

\section{Sintesi}
Dopo una battuta sul Lab~1 (una guida passo passo per imparare a usare una rete, più che una vera
pipeline di data science), si entra nelle architetture deep vere e proprie con le
\textbf{recurrent neural network}. Una FFNN tratta ogni punto in modo indipendente; una RNN mantiene
uno \textbf{stato} che evolve nel tempo e riassume la serie. La lezione costruisce la RNN a partire
dalla FFNN, introduce la \en{backpropagation through time} e i suoi problemi di gradiente, motiva e
descrive l'LSTM, poi passa alla preparazione dei dati (tensore 3D, lunghezze diverse) e alla
progettazione. Approfondimento nel capitolo~\ref{cap:rnn}.

\section{Avvisi pratici}
\begin{itemize}
  \item Il Lab~2 si concentra sul caso \textbf{many-to-one} (una serie in input, una label in output).
  \item L'esempio di codice sulle sequenze prosegue mercoledì (\rif{2026-10-14}{14/10}).
\end{itemize}

\section{Contenuti}

\subsection{Perché le RNN}
Da un solo punto non si può predire la traiettoria: manca la storia. Le RNN gestiscono input
sequenziali di lunghezza variabile e mantengono uno \textbf{hidden state} che rappresenta la serie
fino a $t-1$ (``\en{my brother plays}\ldots'': il contesto suggerisce cosa viene dopo; in modo
semplificatissimo è l'idea della predizione della parola successiva). Applicazioni: traduzione,
previsione di serie (scostamenti forti dalla previsione possono indicare un'anomalia), classificazione
di serie, riconoscimento vocale, \en{part-of-speech tagging}.

\subsection{Varianti e architetture}
Quattro configurazioni: \textbf{one-to-one}, \textbf{many-to-one} (sentiment, sequenza di log per
rilevare un attacco), \textbf{one-to-many} (image captioning), \textbf{many-to-many} (traduzione).
La RNN \textbf{mono-direzionale} va bene per il tempo reale ma fatica con le dipendenze lunghe; la
\textbf{bidirezionale} usa due RNN con pesi diversi in direzioni opposte, cattura passato e futuro ma
costa di più e richiede la sequenza completa (es.\ un'email da classificare come phishing).

\subsection{Dentro la cella}
\[
h_t = F_1\big(W_{hh}h_{t-1} + W_{xh}x_t + b_{hh}\big), \qquad
\hat y_t = F_2\big(W_{yh}h_t + b_{yh}\big)
\]
con $F_1$ tipicamente \textbf{tanh}. Il passaggio dalla FFNN alla RNN: la stessa rete applicata a ogni
istante, poi collegata propagando lo stato. Relazione di ricorrenza $h_t = f_W(x_t, h_{t-1})$:
\textbf{stessa funzione e stessi pesi} a ogni passo. Nel many-to-one solo l'ultima cella ha la
classification head.

\begin{esame}
La parte ricorrente fa \textbf{feature extraction}: lo stato finale comprime la serie in pochi numeri
(un \textbf{embedding} della serie), e la head (da un layer lineare a una FFNN) lo usa per
classificare. Le due parti si addestrano \textbf{insieme}. Il docente insiste anche sul capire
\textbf{perché} la BPTT è più complessa della backpropagation classica.
\end{esame}

\subsection{BPTT e problemi di gradiente}
La \en{backpropagation through time} somma le loss dei passi, propaga l'errore attraverso output e
stati fino all'inizio della sequenza e accumula i gradienti sui pesi condivisi. Il gradiente rispetto
agli stati iniziali contiene molti fattori di $W_{hh}$: valori $>1$ portano all'\textbf{exploding
gradient} (rimedio: \en{gradient clipping}), valori $<1$ al \textbf{vanishing gradient} (rimedi:
attivazione, inizializzazione, architettura). Esempio delle dipendenze lunghe: ``\en{the clouds are in
the}\ldots'' si risolve con il passato recente; una lunga storia che parte dalla Francia e finisce con
``parlo fluentemente\ldots'' no.

\subsection{LSTM}
Proposta nel 1997, mantiene un \textbf{cell state} $C_t$ (memoria a lungo termine) e un
\textbf{hidden state} $h_t$ (breve termine), regolati da gate con sigmoid: \textbf{forget gate}
(quanto passato tenere), \textbf{input gate} e memoria candidata con tanh (quanto presente
aggiungere: un articolo come ``\en{the}'' conta poco), \textbf{output gate} (quanto del cell state
diventa $h_t$). Gestisce dipendenze lunghe al prezzo di molti più pesi. Quanto passato ricordi non si
sa a priori: dipende da dati, training e dimensioni degli stati. I transformer saranno trattati nella
parte di NLP.

\subsection{Dati e progettazione}
L'input di una RNN è un tensore 3D \code{(batch, seq\_len, features)}. Una serie mono-dimensionale si
scrive \code{[[1],[2],[3]]}: \code{[1,2,3]} sarebbe un solo istante con 3 feature. Il numero di
feature per istante è fisso; il numero di passi può variare tra serie, ma dentro un mini-batch serve
la stessa lunghezza (padding o troncamento, solo dentro il batch). Checklist: tipo di RNN, lunghezza
delle sequenze, memoria lunga o no, mono o bidirezionale; input = numero di \textbf{feature}; hidden
size in potenze di 2; 1--2 layer ricorrenti; head lineare o FFNN; ottimizzazione come per le FFNN
(Adam, learning rate 0,001 con scheduler, shuffle solo tra serie indipendenti).

\subsection{Da tabella a sequenze}
Nel dataset sull'inquinamento ogni riga è un'ora (35.040 righe, 7 feature): non ci sono ancora serie.
Una funzione \code{create\_sequences} prende finestre di 7 righe consecutive e come label il valore
di inquinamento successivo, producendo il tensore 3D per il training. L'esempio completo si vede nella
lezione del \rif{2026-10-14}{14/10}.

\section{Domande aperte e punti da verificare}
\begin{verifica}
\begin{itemize}
  \item Label della finestra: a voce il docente dice che è quella dell'``ultimo elemento, il
    settimo''; nel notebook è \code{labels[i + seq\_length]}, cioè il valore \textbf{successivo} alla
    finestra (si predice l'ora dopo). Vale il notebook.
\end{itemize}
\end{verifica}
